\section{API 与 Web Agent 的角色边界}
\subsection{API Agent：明确规范即可自动化}
在 00:08:35--00:09:35 的段落中，讲者把 API Agent 描述为“具名 API 或端点的集合，主要交换文本信息”，其工具集可用公式化规格精确描述，适用于已有服务接口的流程。

\begin{knowledgebox}{API Agent 的核心收益}
当工具集、输入/输出格式明确且可校验时，API Agents 可以保持“高精度、可预测”的执行，同时保持 prompt-engineering 可复现。
\end{knowledgebox}

\subsection{Web Agent：需要感知与 UI 推理}
Web Agents 被描述为“像人一样在浏览器中点击、填写表单、跨页面导航”，但同时承担理解目标、现场感知与长时规划的负担（00:32:49--00:34:14）。

\begin{warningbox}{Web Agent 的三重挑战}
它要理解用户意图、识别当前 UI、还要规划到达最终结果；但终点往往不在当前页面，单凭当前观察无法判定是否已完成，执行错误会导致重复操作或流程中断。
\end{warningbox}

\subsection{本章小结}
API 与 Web Agents 并非互斥，而是互补：前者靠规范保证稳定，后者靠感知实现通用，设计时需要根据目标任务合理组合。

